\section*{अध्याय \ref{ch:b2:genome-diversification} --- उत्परिवर्तन और जीनोम की विविधता}

\begin{solution}{exo:b2:genome-diversification:1}
GAG: मूक (अब भी ग्लूटामेट), एक संक्रमण (A$\to$G)। GTA: कुअर्थी
(वैलीन), एक पर्यावर्तन (A$\to$T)। TAA: निरर्थक (विराम), एक
पर्यावर्तन (G$\to$T)। कोई क्षारक अंतर्वेशित करना: \omterm{def:b2:genome-diversification:mutation}{चौखटा-विस्थापन}, जो
उसके बाद का हर कोडॉन फिर से लिख देता है।
\end{solution}

\begin{solution}{exo:b2:genome-diversification:2}
$U = 5\times 10^{-10}\times 4.6\times 10^{6} = 2.3\times 10^{-3}$:
लगभग 430 में एक विभाजन कोई \omterm{def:b2:genome-diversification:mutation}{उत्परिवर्तन} बनाता है। $10^{9}$ कोशिकाओं में:
प्रति पीढ़ी $2.3\times 10^{6}$ नए \omterm{def:b2:genome-diversification:mutation}{उत्परिवर्तन}।
\end{solution}

\begin{solution}{exo:b2:genome-diversification:3}
\omterm{def:b2:genome-diversification:hgt}{रूपांतरण}: फटी हुई कोशिकाओं से मुक्त DNA का ग्रहण, जिसके लिए कोई सक्षम
ग्राही चाहिए; यह गुणसूत्रीय खंड हस्तांतरित करता है (न्यूमोकॉकस में संपुट
के जीन)। संयुग्मन: किसी पिलस से कोशिका-से-कोशिका संपर्क, जिसके लिए दाता
में कोई संयुग्मी \omterm{def:b2:unicellular-diversity:prokaryote}{प्लास्मिड} चाहिए; यह \omterm{def:b2:unicellular-diversity:prokaryote}{प्लास्मिड}
(प्रतिरोध के जीन) और, किसी Hfr से, गुणसूत्रीय जीन हस्तांतरित करता है। \omterm{def:b2:genome-diversification:hgt}{पारक्रमण}:
कोई जीवाणुभोजी जिसने पोषी DNA संवेष्टित कर लिया हो; यह कोई यादृच्छिक खंड
(सामान्यीकृत) या प्रोफ़ाज स्थल के पड़ोस के जीन (विशिष्टीकृत) ले जाता है।
\end{solution}

\begin{solution}{exo:b2:genome-diversification:4}
\omterm{def:b2:genome-diversification:polyploidy}{समबहुगुणित}: एक ही जाति से अतिरिक्त समुच्चय, किसी द्विगुणित युग्मक के
रास्ते (चतुर्गुणित आलू, अनेक समचतुर्गुणित तिपतिया)। \omterm{def:b2:genome-diversification:polyploidy}{विषमबहुगुणित}: दो
जातियों के समुच्चय, संकरण से जुड़े और फिर दुगुने किए हुए
(रोटी का गेहूँ, कपास, तंबाकू)। द्विगुणित संकर AB में कोई दो गुणसूत्र एक
जैसे नहीं, इसलिए अर्धसूत्री विभाजन I में कोई भी युग्मित नहीं हो सकता;
गुणसूत्र यादृच्छिक रूप से पृथक होते हैं और युग्मक असंतुलित तथा अजीवनक्षम
रहते हैं।
\end{solution}

\begin{solution}{exo:b2:genome-diversification:5}
$2\times 3.2\times 10^{9}\times 1.2\times 10^{-8} = 77$ नए \omterm{def:b2:genome-diversification:mutation}{उत्परिवर्तन}।
कूटलेखी अनुक्रम में: $77\times 0.015 = 1.2$; ऐमीनो अम्ल बदलने वाले:
$1.2\times 0.7 = 0.8$ --- अधिकांश बच्चे लगभग एक नया ऐमीनो-अम्ल
प्रतिस्थापन ढोते हैं।
\end{solution}

\begin{solution}{exo:b2:genome-diversification:6}
घटनाएँ: $m(N - N_0) = 2\times 10^{-8}\times 10^{9} = 20$। उत्परिवर्ती
कोशिकाएँ: $mN\ln(N/N_0) = 20\times \ln 10^{6} = 20\times 13.8 = 276$।
$p_0 = e^{-20} \approx 2\times 10^{-9}$: हर संवर्धन में उत्परिवर्ती हैं,
इसलिए बिना किसी उत्परिवर्ती वाले संवर्धन गिनने से कुछ नहीं मापा जाता।
$m(N -
N_0) \approx 1$ चाहिए, यानी $N \approx 5\times 10^{7}$, जहाँ $p_0 \approx
0.37$ है और कुछ दर्जन संवर्धन $m$ को दो के गुणक के भीतर दे देते हैं।
\end{solution}

\begin{solution}{exo:b2:genome-diversification:7}
यूरेसिल का DNA में कोई स्थान नहीं, इसलिए कोई यूरेसिल-DNA ग्लाइकोसिलेज़
जो भी यूरेसिल मिले उसे बिना किसी द्विविधा के हटा सकता है। थाइमीन एक
सामान्य क्षारक है: 5-मेथिलसाइटोसीन के विऐमीनित होने पर G$\cdot$T बेमेल
दिख तो जाता है, पर मरम्मत तंत्र के पास यह जानने का कोई उपाय नहीं कि
कौन-सा रज्जु ग़लत है, और आधी बार वह G को ही ``सुधार'' देता है और यों
C$\to$T संक्रमण पक्का कर देता है। इसलिए मेथिलित CpG स्थल बाक़ी स्थलों से
लगभग दस गुना तेज़ उत्परिवर्तित होते हैं, और विकास-काल में वे TpG तथा CpA
में बदलते रहे हैं: स्तनधारी जीनोमों में क्षारक-संघटन से अपेक्षित CpG का
केवल पाँचवाँ भाग बचा है।
\end{solution}

\begin{solution}{exo:b2:genome-diversification:8}
प्रति मिनट \qty{46}{kb}। स्थितियाँ: 368, 782, 1150 और
\qty{1840}{kb}; दूरियाँ 414, 368 और \qty{690}{kb}।
\end{solution}

\begin{solution}{exo:b2:genome-diversification:9}
एक गुणसूत्र की दूसरी प्रति दूसरे गुणसूत्र की पहली प्रति से युग्मित होती
है; उनके बीच का जीन-विनिमय एक ऐसी क्रोमैटिड देता है जो प्रति~1, प्रति~2
और प्रति~2$'$ (तीन जीन) ढोती है और एक ऐसी जो केवल
प्रति~1$'$ ढोती है। जिस व्यक्ति को दोनों जनकों से एक-एक जीन वाला
गुणसूत्र मिले उसके पास चार के बजाय दो $\alpha$ जीन हैं
($-\alpha/-\alpha$): $\alpha$ शृंखलाएँ आधी दर से बनती हैं, लाल कोशिकाएँ
छोटी और फीकी होती हैं, और रक्ताल्पता हलकी या नदारद --- यानी
$\alpha$-थैलेसीमिया लक्षण।
\end{solution}

\begin{solution}{exo:b2:genome-diversification:10}
युग्मक: $n = 21$ (A, B, D में से हर एक का एक समुच्चय)। संकर: AB, 14
गुणसूत्र; ABD, 21। रोटी का गेहूँ $\times$ एमर: AABBD, 35
गुणसूत्र --- A और B समुच्चय युग्मित हो जाते हैं, D समुच्चय का कोई साथी
नहीं, और वह पौधा बहुत हद तक बंध्य है। हर दुगुनेपन ने हर गुणसूत्र को एक
समरूप समजात दे दिया, इसलिए द्विसंयोजी बनते हैं और युग्मक संतुलित रहते
हैं; किसी जनक से पीछे संकरण करने पर ऐसे पौधे बनते हैं जिनका एक समुच्चय
अयुग्मित रहता है और जिनके युग्मक असंतुलित होते हैं, इसलिए वह षड्गुणित
अपने पूर्वजों से जननिक रूप से पृथक है --- एक ही पग में एक नई जाति।
\end{solution}

\begin{solution}{exo:b2:genome-diversification:11}
संभार-तांत्रिक वृद्धि: $R(t) = N/\bigl(1 + (N - 1)e^{-\gamma Nt}\bigr)$।
आधी आबादी तब, जब $(N - 1)e^{-\gamma Nt} = 1$: $t = \ln(N -
1)/\gamma N = 20.7/0.5 = \qty{41}{h}$। प्रतिजैविक के बिना
प्लास्मिड-रहित कोशिकाएँ \qty{5}{\%} तेज़ बढ़ती हैं, इसलिए वाहकों और
अवाहकों का अनुपात $e^{-st}$ की तरह क्षीण होता है, जिसमें $s = 0.05\ln 2 =
\qty{0.035}{h^{-1}}$ प्रति घंटा है (\omterm{def:b2:unicellular-diversity:division}{पीढ़ी काल} \qty{1}{h}):
\qty{99}{\%} से \qty{1}{\%} तक आने में $\ln(99^2)/0.035 = \qty{260}{h}$,
यानी ग्यारह दिन लगते हैं --- और लगातार हस्तांतरण उस \omterm{def:b2:unicellular-diversity:prokaryote}{प्लास्मिड}
को किसी कम बारंबारता पर अनिश्चित काल तक बनाए रख सकता है, ताकि वह औषधि
के अगले दौर के लिए तैयार रहे।
\end{solution}

\begin{solution}{exo:b2:genome-diversification:12}
\omterm{thm:b2:genome-diversification:luria}{उच्चावचन परीक्षण} ने दिखाया कि प्रतिरोध के \omterm{def:b2:genome-diversification:mutation}{उत्परिवर्तन} वरण करने वाले
कारक से पहले और उसके बिना, यादृच्छिक समयों पर उठते हैं, इसलिए
\omterm{def:b2:genome-diversification:mutation}{उत्परिवर्तन} इस बात से अंधा है कि उपयोगी क्या होगा; प्रतिलिपि-बुवाई ने
वही बात वरित कोशिकाओं को सामने लाए बिना ही दिखा दी। विकास यादृच्छिक नहीं
है, क्योंकि वरण इन अंधे विभेदों को उनके परिणामों के अनुसार, लगातार और
संचयी रूप से छाँटता है। यादृच्छिकता जहाँ स्वयं गढ़ी जाती है वह एक ही जगह
है, यानी दर: जिन वंशक्रमों में बहुत अधिक या बहुत कम त्रुटियाँ हों वे
हारते हैं, इसलिए \omterm{thm:b2:genome-diversification:rate}{उत्परिवर्तन दर} एक अनुकूलित समझौता है --- पर हर
\omterm{def:b2:genome-diversification:mutation}{उत्परिवर्तन} क्या करेगा, यह अप्रत्याशित और अचयनित ही रहता है।
\end{solution}

\begin{solution}{pb:b2:genome-diversification:1}
\textbf{1.} योग 119, माध्य $5.95$; वर्गों का माध्य $11489/20 = 574$;
प्रसरण $574 - 35 = 540$।
\textbf{2.} योग 324, माध्य $16.2$; वर्गों का माध्य $5462/20 = 273$;
प्रसरण $273 - 262 = 11$।
\textbf{3.} एक ही संवर्धन के प्रतिदर्श पॉइसों हैं (प्रसरण माध्य के
पास): वे एक ही नियत उत्परिवर्ती-अंश से लिए गए यादृच्छिक प्रतिदर्श हैं।
स्वतंत्र संवर्धनों का प्रसरण माध्य का नब्बे गुना है: उनके उत्परिवर्ती
वृद्धि के दौरान भिन्न-भिन्न यादृच्छिक समयों पर उठे, और \omterm{def:b2:genome-diversification:mutation}{उत्परिवर्तन}
जितना जल्दी हुआ उसका क्लोन उतना ही बड़ा।
\textbf{4.} एक जैकपॉट: पहले कुछ विभाजनों में से किसी एक में हुआ कोई \omterm{def:b2:genome-diversification:mutation}{उत्परिवर्तन},
जिसका क्लोन फिर संवर्धन के साथ बढ़कर सौ कोशिकाओं तक पहुँच गया।
\textbf{5.} बीस में से चौदह: $p_0 = 0.70$; $m(N - N_0) = -\ln 0.70
= 0.36$; प्रति विभाजन $m = 0.36/2\times 10^{8} = \qty{1.8e-9}{}$।
\textbf{6.} $1.8\times 10^{-9}\times 2\times 10^{8}\times \ln(2\times
10^{6}) = 0.36\times 14.5 = 5.2$ उत्परिवर्ती कोशिकाएँ; प्रेक्षित $5.95$
--- यानी बीस संवर्धनों के रव के भीतर सहमति।
\textbf{7.} कोई भी प्रतिदर्श ख़ाली नहीं था (सब जनक संवर्धन के
उत्परिवर्ती साझा करते हैं), इसलिए $p_0 = 0$ कोई आकलन नहीं देता; वे
प्रतिदर्श एक ही संवर्धन का उत्परिवर्ती-अंश मापते हैं, जो इस पर निर्भर है
कि उसके \omterm{def:b2:genome-diversification:mutation}{उत्परिवर्तन} कब हुए।
\textbf{8.} प्रति क्षारक युग्म प्रति प्रतिकृतिकरण $u = m/20 = \qty{9e-11}{}$।
\textbf{9.} $U = 9\times 10^{-11}\times 4.6\times 10^{6} =
\qty{4.1e-4}{}$।
\textbf{10.} $10^{12}$ कोशिकाएँ $\times$ $4.1\times 10^{-4}$ $= 4.1\times
10^{8}$ नए \omterm{def:b2:genome-diversification:mutation}{उत्परिवर्तन}।
\textbf{11.} $3\times 4.6\times 10^{6} = 1.4\times 10^{7}$ संभव
परिवर्तन; हर एक $4.1\times 10^{8}/1.4\times 10^{7} \approx 30$
बार उठता है।
\textbf{12.} हर एक-क्षारक परिवर्तन, और उनमें वह हर एक भी जो किसी
\omterm{def:b2:genome-diversification:mutation}{एक-उत्परिवर्तन} से हारने वाली औषधि का प्रतिरोध देता है, औषधि डाले
जाने से पहले ही दर्जनों कोशिकाओं में उपस्थित है; औषधि केवल उन्हें
चुनती है।
\textbf{13.} एक ही कोशिका में दो स्वतंत्र \omterm{def:b2:genome-diversification:mutation}{उत्परिवर्तन} प्रति विभाजन
$m^2 \approx (2\times 10^{-9})^2 = 4\times 10^{-18}$ की दर से होते हैं:
$10^{12}$ कोशिकाओं में प्रति पीढ़ी $4\times 10^{-6}$ --- यानी दोहरी
प्रतिरोधी कोशिका वस्तुतः कभी पहले से नहीं होती, और इसीलिए
यक्ष्मा का उपचार एक साथ कई औषधियों से किया जाता है।
\textbf{14.} $a = N - 1$ और $E = e^{-\gamma Nt}$ के साथ: $R = N/(1 +
aE)$, $\mathrm{d}R/\mathrm{d}t = \gamma N\cdot NaE/(1 + aE)^2$; और
$\gamma R(N - R) = \gamma\,\frac{N}{1 + aE}\cdot\frac{NaE}{1 + aE}$,
यानी वही; $R(0) = N/(1 + N - 1) = 1$।
\textbf{15.} $(N - 1)E = 1$: $t = \ln(N - 1)/\gamma N = 20.7/0.5 =
\qty{41}{h}$।
\textbf{16.} $1 + aE = 1/0.99$, $aE = 0.0101$: $t = \ln(10^{9}/0.0101)
/0.5 = (20.7 + 4.6)/0.5 = \qty{51}{h}$।
\textbf{17.} $\gamma (N/2)^2 = \gamma N\cdot N/4 = 0.5\times 2.5\times
10^{8} = 1.25\times 10^{8}$ हस्तांतरण प्रति घंटा।
\textbf{18.} वृद्धि दरें $\ln 2 = \qty{0.693}{h^{-1}}$ और
\qty{0.762}{h^{-1}}; उत्परिवर्ती/संवेदी अनुपात $e^{st}$ की तरह बढ़ता है,
जिसमें $s = \qty{0.069}{h^{-1}}$; $1/N$ से 1 तक $\ln N/s =
20.7/0.069 = \qty{300}{h}$, यानी बारह दिन लगते हैं।
\textbf{19.} हस्तांतरण वाहकों और अवाहकों के गुणनफल के समानुपाती है और
मुलाक़ातों की दर से चलता है, विभाजनों की दर से नहीं; वंशानुक्रम
उत्परिवर्ती के वृद्धि-लाभ से सीमित है, जो छोटा है।
कोई \omterm{def:b2:unicellular-diversity:prokaryote}{प्लास्मिड} आबादी दो दिन में पार कर लेता है, कोई अनुकूलित उत्परिवर्ती
दो सप्ताह में।
\textbf{20.} $U = 4.1\times 10^{-2}$; हानिकारक: प्रति पीढ़ी $0.041\times 0.88\times
0.3 = 0.011$।
\textbf{21.} माध्य $2.2$: $e^{-2.2} = 0.11$, यानी नौ में एक संतान
अक्षत।
\textbf{22.} वन्य प्ररूप: प्रति पीढ़ी $1.1\times 10^{-4}$, 200 में $0.022$:
$e^{-0.022} = 0.98$।
\textbf{23.} माध्य $520$ उत्परिवर्ती कोशिकाएँ; $p_0 = e^{-36} \approx
10^{-16}$: कोई ख़ाली संवर्धन नहीं।
\textbf{24.} तीव्र, बदलते वरण के नीचे (किसी अस्पताल के बारी-बारी बदले
जाते प्रतिजैविक) नए विभेदों की सौ गुनी आपूर्ति बोझ से भारी पड़ती है, और
उत्परिवर्तक अपने बनाए प्रतिरोधों के साथ मुफ़्त सवारी कर लेते हैं। किसी
स्थिर पर्यावरण में पाने को कुछ नहीं है और खोने को प्रति सौ पीढ़ी एक
हानिकारक \omterm{def:b2:genome-diversification:mutation}{उत्परिवर्तन} है, इसलिए उत्परिवर्तक वंशक्रम छाँट दिए जाते हैं।
\textbf{25.} प्रतिरोध तक प्रति विभाजन $m = \qty{1.8e-9}{}$, प्रति
क्षारक युग्म प्रति प्रतिकृतिकरण $u =
\qty{9e-11}{}$, प्रति जीनोम प्रति प्रतिकृतिकरण $U = \qty{4.1e-4}{}$
\omterm{def:b2:genome-diversification:mutation}{उत्परिवर्तन}; वह \omterm{def:b2:unicellular-diversity:prokaryote}{प्लास्मिड} आँत के आधे
भाग तक \qty{41}{h} में और उसके \qty{99}{\%} तक \qty{51}{h} में पहुँचता है।
\end{solution}
